% Phase-2 fix designs, group G1 (set theory): R-001, R-017.
% Fragment: compiles inside the gap-analysis document under amsmath/amssymb.

\subsection*{R-001 --- Honest CH independence (v1.0 lines 9062--9133; chapter 8665--9133)}

\paragraph{Defect.}
The theorem at 9065 claims ``the constructions using \(L\) and forcing establish'' \(\mathsf{Con}(\mathsf{ZFC})\Rightarrow\mathsf{Con}(\mathsf{ZFC}+\neg\mathsf{CH})\), but the only forcing presented is over a countable transitive model, which Warning 9107--9117 says ``assumes more than the bare arithmetical statement \(\mathsf{Con}(\mathsf{ZFC})\)'' and must be discharged; it never is.  The sketch (9092--9094) derives nonderivability by soundness from ``the two structures satisfying the displayed theories,'' but neither is a set-sized structure: \(L\) is a proper class and \(M\) is hypothesized (so too the completeness appeal at 8759--8760).

\paragraph{Repair.}
\textbf{(1) 8733--8762.}  Keep ``\(L\) is transitive and contains every ordinal'' (\emph{Status: proof sketch}).  Delete the \(\mathsf{Con}\) displays (8739--8750) and the completeness sentence (8758--8760).  Replace the satisfaction display by a separately stated scheme (\emph{Status: cited theorem}; Jech 2003, Ch.~13, pp.~175--200 --- reflection, definable well-order, condensation, as \(\mathsf{ZF}\)-theorems about the definable \(L\)): \(\mathsf{ZF}\vdash\sigma^{L}\) for every axiom \(\sigma\) of \(\mathsf{ZFC}+\mathsf{GCH}+(V{=}L)\), with \((\cdot)^{L}\) the convention of 8710--8731.
\textbf{(2) 8764--8782.}  Retitle the section ``Forcing over a countable transitive set: a pedagogical device.''  After the convention insert a Remark (\emph{Status: proof sketch}; omitted verifications: the translation convention, the absoluteness computation, and the interpreted derivability-condition bookkeeping).  \(\mathsf{Con}(\mathsf{ZFC})\) (as at 7808) is a sentence of arithmetic, and the manuscript fixes no translation of arithmetic into the membership language, so the Remark must first state one: numbers as finite ordinals under the canonical interpretation, writing \(\sigma^{\in}\) for the membership-language translation of an arithmetic sentence \(\sigma\).  Route: if such an \(M\) exists, soundness for the set-sized \((M,\in)\) gives metatheoretic consistency of \(\mathsf{ZFC}\); correctness of the fixed proof predicate (Warning at 7818--7824) makes \(\mathsf{Con}(\mathsf{ZFC})\) true in the standard structure; absoluteness of arithmetic over transitive models containing \(\omega\) (Jech, Ch.~12, pp.~155--174) then gives \((M,\in)\models\mathsf{ZFC}+\mathsf{Con}(\mathsf{ZFC})^{\in}\), and soundness again gives consistency of \(\mathsf{ZFC}+\mathsf{Con}(\mathsf{ZFC})^{\in}\).  That consistency sentence is unprovable in \(\mathsf{ZFC}+\mathsf{Con}(\mathsf{ZFC})^{\in}\) itself --- not by \texttt{thm:godel-second} verbatim, which covers computably enumerable \(T\supseteq\mathsf Q\) in the language of arithmetic and so does not apply to a membership-language theory, but by that theorem transported across the canonical interpretation, under which the fixed provability predicate still satisfies the derivability conditions (this bookkeeping is the third omitted verification); so the hypothesis strictly exceeds \(\mathsf{Con}(\mathsf{ZFC})\); other section results stay unchanged.
\textbf{(3) 9062--9096.}  Delete the theorem and sketch (9065--9096); insert:
(a) \emph{Lemma (Relativization of derivations)} [\emph{complete proof}]: for membership-language formulas \(\varphi_1,\dots,\varphi_k,\psi\) with free variables among \(v_1,\dots,v_m\), if \(\{\varphi_1,\dots,\varphi_k\}\vdash\psi\) in the fixed natural-deduction calculus, then \(\mathsf{ZF}\vdash\forall v_1\cdots\forall v_m\,(\bigwedge_{j\le m} v_j\in L\wedge\bigwedge_{i\le k}\varphi_i^{L}\to\psi^{L})\).
Route (induction on the derivation): \((\cdot)^{L}\) commutes with connectives and \(\bot^{L}=\bot\), so propositional and classical rules lift verbatim; terms are variables, so \(\forall\)-elimination and \(\exists\)-introduction consume a guard \(u\in L\) from the antecedent, while \(\forall\)-introduction and \(\exists\)-elimination form the guarded clauses of 8718--8723 (eigenvariable condition unchanged; \(x\in L\) is parameter-free); a free variable vanishing from the conclusion is discharged via \(\mathsf{ZF}\vdash\exists u\,(u\in L)\), witness \(\varnothing\in L_1\) (justify inside \(\mathsf{ZF}\) by the comprehension reading of \(\operatorname{Def}\): \(\{x\in L_0:\ldots\}=\varnothing\) by separation whatever the satisfaction clause returns on the empty domain excluded at v1.0 3028; do not route through the proof-sketch result that \(L\) contains every ordinal); equality rules act on unchanged atoms; the explicit \(\alpha\) rule (v1.0 3475--3487) maps to an \(\alpha\) step, since relativization inserts the same parameter-free guard under alpha-variant binders.  For \(m=0\) or \(k=0\) the corresponding conjunction is absent from the displayed conclusion.
(b) \emph{Theorem (Relative consistency of \(\mathsf{CH}\))} [\emph{complete proof}]: if \(\mathsf{ZF}\) is consistent, so is \(\mathsf{ZFC}+\mathsf{GCH}+(V{=}L)\); hence if \(\mathsf{ZFC}\) is consistent, so is \(\mathsf{ZFC}+\mathsf{CH}\).  Proof: a derivation of \(\bot\) from axioms \(\sigma_1,\dots,\sigma_k\), the Lemma (\(m=0\), \(\psi=\bot\)), and the item-1 scheme give \(\mathsf{ZF}\vdash\bot\); \(\mathsf{ZF}\)-derivations are \(\mathsf{ZFC}\)-derivations, and \(\mathsf{ZFC}+\mathsf{GCH}+(V{=}L)\) derives \(\mathsf{CH}\) (the \(\alpha=0\) instance of 8633), so inconsistency of \(\mathsf{ZFC}+\mathsf{CH}\) lifts to it.
(c) \emph{Cited theorem (Relative consistency of \(\neg\mathsf{CH}\))} [\emph{cited theorem}; Jech, Ch.~14, pp.~201--224]: if \(\mathsf{ZFC}\) is consistent, so is \(\mathsf{ZFC}+\neg\mathsf{CH}\).  Imported as \(\mathsf{ZFC}\)-theorem schemes about the definable class \(V^{B}\), \(B\) the completion of the poset of finite partial functions \(\omega_2\times\omega\to 2\): (i) \(\lVert\sigma\rVert_B=1\) for each \(\mathsf{ZFC}\) axiom; (ii) Boolean soundness: \(\{\sigma_1,\dots,\sigma_k\}\vdash\psi\) implies \(\mathsf{ZFC}\vdash\lVert\sigma_1\rVert\wedge\cdots\wedge\lVert\sigma_k\rVert\le\lVert\psi\rVert\), with \(\lVert\bot\rVert=0\) and \(B\) provably nontrivial; (iii) \(B\) is ccc and ccc forcing preserves cardinals: \(\lVert\check\aleph_1=\aleph_1\wedge\check\aleph_2=\aleph_2\rVert=1\); (iv) the distinctness density argument (9048--9055, per \(\omega_2\)-section) gives \(\lVert 2^{\aleph_0}\ge\check\aleph_2\rVert=1\), hence \(\lVert\neg\mathsf{CH}\rVert_B=1\).  A \(\bot\)-derivation from \(\mathsf{ZFC}\cup\{\neg\mathsf{CH}\}\) then yields \(\mathsf{ZFC}\vdash 1=0\) in \(B\): contradiction.  This is the ``Boolean-valued'' discharge Warning 9107 demands --- no countable transitive model; consistency is calculus-invariant via soundness and strong completeness.
(d) \emph{Remark (Arithmetized form)} [\emph{proof sketch}]: the two derivation transformations are primitive recursive, so \(\mathsf{Con}(\mathsf{ZFC})\to\mathsf{Con}(\mathsf{ZFC}+\mathsf{CH})\) and \(\mathsf{Con}(\mathsf{ZFC})\to\mathsf{Con}(\mathsf{ZFC}+\neg\mathsf{CH})\) are \(\mathsf{PA}\)-theorems; omitted: the arithmetization, including a primitive recursive translation between the fixed natural-deduction calculus and the source calculus of (c) --- the semantic calculus-invariance argument in (c) does not supply this; the deduction-theorem translation does.
(e) \emph{Theorem (Nonderivability)} [\emph{complete proof}]: if \(\mathsf{ZFC}\) is consistent then \(\mathsf{ZFC}\nvdash\mathsf{CH}\) and \(\mathsf{ZFC}\nvdash\neg\mathsf{CH}\).  Proof: \(\mathsf{ZFC}\vdash\mathsf{CH}\) with monotonicity and the assumption rule gives \(\mathsf{ZFC}+\neg\mathsf{CH}\vdash\mathsf{CH},\neg\mathsf{CH}\), negation elimination gives \(\vdash\bot\), and the contrapositive of (c) refutes consistency of \(\mathsf{ZFC}\); symmetrically \(\mathsf{ZFC}\vdash\neg\mathsf{CH}\) makes \(\mathsf{ZFC}+\mathsf{CH}\vdash\bot\), contradicting (b).
\textbf{(4) 9098--9117.}  Retain both Warnings, appending to the second that (c) supplies the Boolean-valued discharge; Exercise 9126 and the sketch at 13941--13952 remain correct.

\paragraph{Acceptance.}
(1) The section replacing 9065--9096 contains exactly the five results of item (3) with the stated statuses; the phrases ``the constructions using \(L\) and forcing establish'' and ``the two structures satisfying the displayed theories'' and the completeness sentence at 8759--8760 no longer occur in 8665--9133.
(2) No proof in the exact-status section uses the countable-transitive-model convention, the generic-filter lemma, or \(L\models\); nonderivability uses only monotonicity, assumption, negation elimination, and contrapositives of (b), (c).
(3) Jech citations carry pages 175--200, 201--224, 155--174; the retitled document builds cleanly.

\subsection*{R-017 --- Split the five-clause finite-coding theorem by status (v1.0 lines 678--702)}

\paragraph{Defect.}
The five-clause theorem \texttt{thm:infinite-cardinal-finite-coding} carries the single label ``Status: cited theorem,'' but the source is invoked only ``for complete proofs of the cardinal products and union bounds'' (clauses 1--4); clause 5 rests on the unlabeled inline sentence ``The final item follows by assigning each generated expression its finite syntax string and applying the preceding items'' (699--700).  A result the manuscript itself proves is hidden inside a cited status, violating the policy at 313--319 that each result carries exactly one of the three status terms and that in-document dependencies be stated separately.

\paragraph{Repair.}
\begin{enumerate}
\item Restate the theorem with clauses 1--4 only, statements unchanged, label kept, ``\emph{Status: cited theorem}''; cite Jech 2003, Theorem~3.5 (\(\aleph_\alpha\cdot\aleph_\alpha=\aleph_\alpha\); with the already-cited well-ordering theorem every infinite cardinal is an aleph) and Ch.~5 (cardinal sums and the union bound \(\lvert\bigcup_{i\in I}A_i\rvert\le\lvert I\rvert\cdot\sup_{i\in I}\lvert A_i\rvert\) under \(\mathsf{AC}\)).  The postamble 698--702 shrinks to this citation sentence plus the retained no-infinite-products caveat; the inline clause-5 sentence is deleted.
\item Immediately after (landing at v1.0 line 697) add \texttt{cor:syntax-counting}, ``\emph{Status: complete proof}'': \emph{Assume the Axiom of Choice and let \(\kappa\) be an infinite cardinal.  (i) Every set \(A\) with \(\lvert A\rvert\le\kappa\) satisfies \(\lvert A^{<\omega}\rvert\le\kappa\), where \(A^{<\omega}\) is the set of finite sequences from \(A\).  (ii) If \(\Sigma\) is an alphabet with \(\lvert\Sigma\rvert\le\kappa\), then the finite strings over \(\Sigma\), the finite tuples of such strings, and the formulas generated by any fixed finitary grammar whose alphabet extends \(\Sigma\) by at most countably many further symbols each form a set of cardinality at most \(\kappa\).}
Proof (promoted from the inline argument): (i) fix an injection \(f\colon A\to\kappa\) (the meaning of \(\lvert A\rvert\le\kappa\)); then \((a_1,\dots,a_n)\mapsto(f(a_1),\dots,f(a_n))\) is a length-preserving injection of \(A^{<\omega}\) into \(\kappa^{<\omega}\), and clause 2 gives \(\lvert\kappa^{<\omega}\rvert=\kappa\).  (ii) Strings over \(\Sigma\) lie in \(\Sigma^{<\omega}\): apply (i) with \(A=\Sigma\).  Tuples of strings lie in \((\Sigma^{<\omega})^{<\omega}\): apply (i) with \(A=\Sigma^{<\omega}\), of cardinality \(\le\kappa\) by the previous step.  Every generated expression is a finite string over the grammar's full alphabet (the manuscript's own fact, v1.0 line 1126); writing \(X\) for the at most countable set of adjoined symbols hypothesized in the statement (logical symbols, punctuation, and variables), clause 4 (applied to \(\Sigma,X,\varnothing,\varnothing,\dots\)) gives \(\lvert\Sigma\cup X\rvert\le\kappa\), so the formula set injects into \((\Sigma\cup X)^{<\omega}\) and (i) applies.
\item Retarget the clause-5 uses: v1.0 lines 3949 (Henkin--Lindenbaum, ``at most \(\kappa\) formulas'') and 4347 (downward L\"owenheim--Skolem, ``at most \(\mu\) function symbols, formulas, and finite tuples'') now cite \texttt{cor:syntax-counting}; line 3969 (clause 3) keeps its explicit reference to \texttt{thm:infinite-cardinal-finite-coding}; at 4350, the anaphor ``The same theorem gives \(\lvert N\rvert\le\mu\)'' (clause 4) currently resolves through the citation at 4347 and would silently retarget with it, so it must be rewritten as an explicit cross-reference to \texttt{thm:infinite-cardinal-finite-coding}.
\end{enumerate}

\paragraph{Acceptance.}
(1) The region replacing 678--702 contains exactly two labeled results: the four-clause theorem (``Status: cited theorem,'' Jech citation naming the product and union-bound facts) and \texttt{cor:syntax-counting} (``Status: complete proof'' with a proof environment containing the injection argument); the sentence beginning ``The final item follows'' no longer occurs.
(2) The corollary's proof appeals only to clauses 2 and 4, the fixed injection, and the formulas-are-strings fact, with no citation to Jech inside the proof.
(3) References formerly at 3949 and 4347 resolve to \texttt{cor:syntax-counting}, those at 3969 and 4350 to \texttt{thm:infinite-cardinal-finite-coding}; the document compiles with no undefined references.
